%% kapa: exam=1 round=21 year=2010 period=1 form=A subject=economics subject_ko=경제원론 seq=3 number=43
%% source: Q-Net 감정평가사 기출문제 (cst00309) artlSeq=5024181 file=extracted/1008/5024181/2010년 제21회 감정평가사 1차시험 기출문제/1교시.hwp
%% conversion: kordoc 4.24.51 -> tools/gen1.py (auto); review=proofread-sonnet; flags=none
\begin{question}{43}
폭설로 도로가 막혀 교통이 두절되고 농촌 비닐하우스가 무너져 농작물 피해가 발생하였다. 우하향하는 총수요곡선과 우상향하는 총공급곡선을 이용하여 이러한 자연재해가 단기적으로 경제에 미치는 영향은?
\begin{choices}
  \item 물가수준은 상승하고 실질 GDP는 감소한다.
  \item 물가수준은 하락하고 실질 GDP는 감소한다.
  \item 물가수준은 상승하고 실질 GDP는 증가한다.
  \item 물가수준은 상승하고 실질 GDP는 불변이다.
  \item 물가수준은 하락하고 실질 GDP는 증가한다.
\end{choices}
\end{question}
\answer{1}
